% ===== PRÉAMBULE CANONIQUE — à coller VERBATIM en tête de CHAQUE .tex =====
\documentclass[11pt,a4paper]{article}
\usepackage[utf8]{inputenc}\usepackage[T1]{fontenc}\usepackage{lmodern}
\usepackage[french]{babel}
\usepackage{amsmath,amssymb,mathtools}\usepackage{icomma}
\usepackage{siunitx}\sisetup{locale=FR}  % \qty{}{}, \si{}, \ang{}
\usepackage{xcolor}\usepackage{colortbl}
\usepackage{tikz}\usepackage{pgfplots}\pgfplotsset{compat=1.18}
\usepackage[siunitx,europeanresistors]{circuitikz} % circuits (élec.)
\usepackage[most]{tcolorbox}\usepackage{enumitem}\usepackage{array,booktabs}
\usepackage[a4paper,margin=2.2cm]{geometry}
\usetikzlibrary{arrows.meta,calc,angles,quotes,positioning,patterns,%
  backgrounds,shapes.geometric,decorations.pathmorphing,decorations.markings,intersections}
\definecolor{primary}{RGB}{222,49,99}\definecolor{primarydark}{RGB}{150,22,55}
\definecolor{defcol}{RGB}{0,114,178}\definecolor{methcol}{RGB}{52,92,148}
\definecolor{excol}{RGB}{0,135,90}\definecolor{attncol}{RGB}{200,40,55}
\tcbset{boxbase/.style={enhanced,breakable,boxrule=0.9pt,arc=2.5pt,
  left=3mm,right=3mm,top=2.3mm,bottom=2.3mm,fonttitle=\bfseries,coltitle=white}}
% --- boîtes TUTORIEL (noms EXACTS, ne pas renommer) ---
\newtcolorbox{cle}[1][]{boxbase,colback=defcol!6,colframe=defcol,
  title={\textsc{À retenir}\ifx\relax#1\relax\relax\else\textmd{ : \itshape #1}\fi}}
\newtcolorbox{methode}[1][]{boxbase,colback=methcol!7,colframe=methcol,sharp corners,
  title={\textsc{Méthode}\ifx\relax#1\relax\relax\else\textmd{ --- \itshape #1}\fi}}
\newtcolorbox{exemplebox}[1][]{boxbase,colback=excol!5,colframe=excol,
  title={\textsc{Exemple}\ifx\relax#1\relax\relax\else\textmd{ : \itshape #1}\fi}}
\newtcolorbox{piege}[1][]{boxbase,colback=attncol!6,colframe=attncol,
  title={\textsc{Piège}\ifx\relax#1\relax\relax\else\textmd{ : \itshape #1}\fi}}
\newtcolorbox{remarque}[1][]{boxbase,colback=black!4,colframe=black!45,
  title={\textsc{Remarque}\ifx\relax#1\relax\relax\else\textmd{ : \itshape #1}\fi}}
% --- boîtes EXERCICE / CORRIGÉ (noms EXACTS, auto-comptées) ---
\newtcolorbox[auto counter]{exo}[1][]{boxbase,colback=primary!4,colframe=primary,
  title={\textsc{Exercice}~\thetcbcounter\ifx\relax#1\relax\relax\else\textmd{ --- \itshape #1}\fi}}
\newtcolorbox[auto counter]{corrige}[1][]{boxbase,colback=excol!4,colframe=excol,
  title={\textsc{Corrigé}~\thetcbcounter\ifx\relax#1\relax\relax\else\textmd{ --- \itshape #1}\fi}}
\newcommand{\vv}[1]{\vec{#1}}\newcommand{\ds}{\displaystyle}
\newcommand{\R}{\mathbb{R}}\newcommand{\N}{\mathbb{N}}\newcommand{\Z}{\mathbb{Z}}
% (un \usepackage{fancyhdr}+en-tête est possible ; il est ignoré côté web)
% =========================================================================
\begin{document}
% THEME_TITLE: L'énergie cinétique et son théorème
\sisetup{per-mode=symbol}

L'\textbf{énergie cinétique} est l'énergie que possède un corps \emph{du fait de son mouvement}. Ce
thème installe sa formule, le \textbf{théorème de l'énergie cinétique} (qui relie travail et
variation de vitesse) et sa dépendance \emph{au carré} de la vitesse.

\begin{cle}[énergie cinétique]
\[
\boxed{\;E_c=\tfrac12\,m\,v^2\;}\qquad(\text{en }\si{\joule})
\]
avec $m$ en \si{\kilogram} et $v$ en \si{\metre\per\second}. L'énergie cinétique est \emph{toujours
positive} et nulle au repos. \textbf{Attention} : $v$ doit être en \si{\metre\per\second} (diviser les
\si{\kilo\metre\per\hour} par $3,6$).
\end{cle}

\begin{cle}[théorème de l'énergie cinétique]
Dans un référentiel galiléen, la somme des travaux de \emph{toutes} les forces entre $A$ et $B$ égale
la variation d'énergie cinétique :
\[
\boxed{\;\sum W_{AB}(\vv F)=E_c(B)-E_c(A)=\tfrac12 m v_B^2-\tfrac12 m v_A^2\;}
\]
Si $\sum W>0$ le corps \emph{accélère} ; si $\sum W<0$ il \emph{ralentit} ; si $\sum W=0$ sa vitesse
est conservée. C'est l'outil qui relie directement les \emph{forces} à la \emph{variation de vitesse},
sans passer par l'accélération.
\end{cle}

\begin{methode}[que faire avec le théorème ?]
\begin{itemize}[leftmargin=1.5em,topsep=1pt,itemsep=1pt]
  \item \textbf{Trouver une vitesse} : on calcule $\sum W$, puis $v_B=\sqrt{v_A^2+\dfrac{2\sum W}{m}}$.
  \item \textbf{Trouver une distance de freinage} : à l'arrêt $v_B=0$, le travail résistant
        $-F\,d$ absorbe toute l'énergie cinétique : $F\,d=\tfrac12 m v_A^2$.
  \item \textbf{Trouver une force de frottement} : $F_f=\dfrac{\tfrac12 m v_A^2}{d}$.
\end{itemize}
\end{methode}

\begin{piege}[la dépendance en $v^2$]
La vitesse intervient \emph{au carré} : si $v$ est multipliée par $n$, alors $E_c$ est multipliée par
$n^2$.
\[
v\to n\,v\ \Longrightarrow\ E_c\to n^2 E_c .
\]
Tripler la vitesse ($n=3$) multiplie $E_c$ par $9$ ; comme la distance de freinage est proportionnelle
à $E_c$, elle est elle aussi multipliée par $9$. C'est pourquoi la distance d'arrêt d'un véhicule
augmente si vite avec la vitesse.
\end{piege}

\begin{center}
\begin{tikzpicture}
  \begin{axis}[width=10cm,height=4.6cm,axis lines=left,
     xlabel={\small vitesse $v$ (\si{\metre\per\second})},ylabel={\small $E_c$ (\si{\joule})},
     xmin=0,xmax=10.5,ymin=0,ymax=55,xtick={0,2,4,6,8,10},ytick={0,10,20,30,40,50},
     tick label style={font=\footnotesize},label style={font=\footnotesize},
     domain=0:10,samples=60]
    \addplot[defcol,line width=1.3pt]{0.5*x^2};
    \node[defcol] at (axis cs:3.9,33){\small $E_c=\tfrac12 m v^2\ (m=\qty{1}{\kilogram})$};
    \draw[dashed,black!45] (axis cs:3,0)--(axis cs:3,4.5)--(axis cs:0,4.5) node[left,black]{\scriptsize $E_c$};
    \draw[dashed,black!45] (axis cs:9,0)--(axis cs:9,40.5)--(axis cs:0,40.5) node[left,black]{\scriptsize $9E_c$};
    \fill[primary] (axis cs:3,4.5) circle (1.7pt); \fill[primary] (axis cs:9,40.5) circle (1.7pt);
  \end{axis}
\end{tikzpicture}
\end{center}

\begin{exemplebox}[augmenter la vitesse d'un chariot]
Un chariot de $m=\qty{10}{\kilogram}$ passe de $v_A=\qty{1}{\metre\per\second}$ à
$v_B=\qty{3}{\metre\per\second}$. Le travail à fournir est
$W=\Delta E_c=\tfrac12\times10\times(3^2-1^2)=\tfrac12\times10\times8=\qty{40}{\joule}$.
\end{exemplebox}
\end{document}
